% !LW recipe = make-deck
% Level 2T (proofs track, TPIL): selection and union, proved. Theory:
% conditions, selection, union, laws 1-3 with proofs. Lean: tables as
% properties of rows (namespace Sets); laws 1-3 as ↔, in term mode.
% Build: make (rules in ../common/deck.mk).
\DocumentMetadata{
  lang        = en-GB,
  pdfstandard = ua-2,
  tagging     = on,
}
\documentclass{ltx-talk}

\input{preamble}

\title{Level 2T: Selection and Union, Proved}
\author{The Proofs and Programs Club}
\institute{Royal Holloway, University of London}

\begin{document}

\leveltitle{2T}{Selection and Union, Proved}{Proofs track: follows TPIL}
  {After \levellink{1}}{Next: \levellink{3t}}

\levelframe{2T}{Level 2T at a Glance}{%
  tables as properties of rows, and selection and union on them.}{%
  laws 1, 2 and 3, for every table.}{%
  read a proof of ``if $A$ then $B$'' as a function from proofs of $A$ to
    proofs of $B$ & \rds{tpil-pat} \\
  prove ``and'', ``or'' and ``if and only if'', and use them &
    \rds{tpil-proplogic} \\
}

% ---------------------------------------------------------------- theory
\theorypart{Selection, union, and three laws}

\begin{frame}
  \frametitle{A Condition Is a Property of Rows}
  \defn{Definition 2T.1 (Condition)}{A \emph{condition} on rows of
    $D_1 \times D_2$ is a property $p$ of rows: for each row $t$, $p(t)$
    is a statement, which holds or does not.}
  For example, on $\rel{Member}$ and on $\rel{Attended}$:
  \[
    \mathit{cs}(n, d) \text{ holds when } d = \text{CS}, \qquad
    \mathit{late}(n, w) \text{ holds when } w \geq 2.
  \]
  \defn{Definition 2T.2 (Both conditions)}{$(p \wedge q)(t)$ holds exactly
    when $p(t)$ holds and $q(t)$ holds.}
  \theorytag{2T}
\end{frame}

\begin{frame}
  \frametitle{Selection Keeps the Rows That Satisfy a Condition}
  \defn{Definition 2T.3 (Selection)}{For a table $R$ and a condition $p$,
    \ $\sel{p}(R) = \{t \in R \mid p(t)\}$.}
  {\centering
    \smalltable{$\sel{\mathit{cs}}(\rel{Member})$}{ll}
      {$\att{name}$ & $\att{dept}$}
      {Alan & CS \\ Grace & CS}\hspace{16mm}
    \smalltable{$\sel{\mathit{late}}(\rel{Attended})$}{lr}
      {$\att{name}$ & $\att{week}$}
      {Alan & 2 \\ Grace & 3}\par}
  \medskip
  A selection never adds rows: $\sel{p}(R) \subseteq R$. In SQL, a
  selection is a \code{WHERE} clause.
  \slideref{\lnk{ahv4}}
  \theorytag{2T}
\end{frame}

\begin{frame}
  \frametitle{Two Selections Are One: Law 1}
  \defn{Law 1}{For every table $R$ and conditions $p$ and $q$:
    \ $\sel{p}(\sel{q}(R)) = \sel{p \wedge q}(R)$.}
  \emph{Proof.} Two sets are equal when they have the same elements. For
  every row $t$:
  \[
    \begin{array}{@{}r@{\quad}l@{}}
      t \in \sel{p}(\sel{q}(R)) & \iff t \in \sel{q}(R) \text{ and }
        p(t) \\
      & \iff t \in R \text{ and } q(t) \text{ and } p(t) \\
      & \iff t \in \sel{p \wedge q}(R). \\
    \end{array}
  \]
  \proofend\par
  Each step unfolds a definition: 2T.3, then 2T.3 again, then 2T.2 and
  2T.3.
  \defn{Law 2}{$\sel{p}(\sel{q}(R)) = \sel{q}(\sel{p}(R))$.
    \textbf{Try}: prove it on paper, in the same way.}
  \theorytag{2T}
\end{frame}

\begin{frame}
  \frametitle{Union Combines Two Tables; Selection Distributes}
  \defn{Definition 2T.4 (Union)}{For tables $R$ and $S$ of the same rows:
    \ $R \cup S = \{t \mid t \in R \text{ or } t \in S\}$.}
  \defn{Law 3}{$\sel{p}(R \cup S) = \sel{p}(R) \cup \sel{p}(S)$.}
  \emph{Proof.}
  \[
    \begin{array}{@{}r@{\quad}l@{}}
      t \in \sel{p}(R \cup S) & \iff (t \in R \text{ or } t \in S)
        \text{ and } p(t) \\
      & \iff (t \in R \text{ and } p(t)) \text{ or } (t \in S
        \text{ and } p(t)) \\
      & \iff t \in \sel{p}(R) \cup \sel{p}(S). \\
    \end{array}
  \]
  \proofend
  \slideref{\lnk{ahv5}}
  \theorytag{2T}
\end{frame}

% ---------------------------------------------------------------- Lean
\leanpart{Tables as properties, and the three laws}

\begin{frame*}
  \frametitle{In Lean, a Table Can Be a Property of Rows}
  \leanblock{sets.lean}
  \begin{itemize}
    \item \code{R t} is the statement ``$t$ is a row of $R$''. So
      \code{select} and \code{union} are formulas, with \code{∧} and
      \code{∨}.
    \item Each name inside \code{namespace Sets} starts with \code{Sets.}:
      put your theorems inside too.
  \end{itemize}
  \leantag{2T}
\end{frame*}

\begin{frame*}
  \frametitle{Your First Proof About Every Table}
  \leanblock{first.lean}
  \begin{itemize}
    \item A proof of ``if $A$ then $B$'' is a function: from a proof
      \code{h} of $A$, it makes a proof of $B$.
    \item A proof \code{h} of \code{A ∧ B} is a pair: \code{h.1} proves
      \code{A}, \code{h.2} proves \code{B}, and \code{⟨a, b⟩} proves
      \code{A ∧ B}.
    \item \ex{2T.1} Prove \code{select p R t → p t}.
  \end{itemize}
  \slideref{\lnk{tpil-pat} \rfsep \lnk{tpil-proplogic}}
  \leantag{2T}
\end{frame*}

\begin{frame*}
  \frametitle{Your Turn: Laws 1 and 2}
  \leanblock{laws12.lean}
  \begin{itemize}
    \item A proof of \code{A ↔ B} is a pair of functions: \code{⟨fun h =>
      \ldots, fun h => \ldots⟩}.
    \item \ex{2T.2} Prove law 1. \ex{2T.3} Prove law 2.
  \end{itemize}
  \slideref{\lnk{tpil-proplogic}}
  \leantag{2T}
\end{frame*}

\begin{frame*}
  \frametitle{Your Turn: Law 3, with ``Or''}
  \leanblock{law3.lean}
  \begin{itemize}
    \item \code{Or.inl a} proves \code{A ∨ B} from a proof \code{a} of
      \code{A}; \code{Or.inr b}, from a proof \code{b} of \code{B}.
    \item To use \code{h : A ∨ B}, take both cases:
      \code{h.elim (fun a => \ldots) (fun b => \ldots)}.
    \item \ex{2T.4} Prove law 3.
  \end{itemize}
  \slideref{\lnk{tpil-proplogic}}
  \leantag{2T}
\end{frame*}

\end{document}
